\section{Orientation tritique, prolongement compatible et double projection}
\label{sec:trit-projection}

Le caractère engendré des coordonnées réelles s’exprime dans la relation
entre profondeur arithmétique et évaluation. La donnée primitive est une histoire
admissible d’opérations avec leurs restes, quotients et orientations. Une
lecture archimédienne associe à ses préfixes des intervalles de plus en plus étroits ;
une lecture incidentielle conserve des relations de frontière du même état.
La double projection du corpus \cite{hmtI,HMT-IX} réunit ces deux opérations
sans convertir une évaluation scalaire en la totalité de l’histoire.
L’article utilise cet antécédent pour étudier les réalisations ultérieures
du registre dodécaphasique.

\subsection{Cylindres positionnels et détermination de la coordonnée réelle}

Pour un bloc ternaire \(b=(b_1,\ldots,b_6)\), avec des représentants
\(b_j\in\{0,1,2\}\) de \(\mathbb F_3\), l’évaluation entière
\(\nu(b)=\sum_{j=1}^6b_j3^{6-j}\) appartient à
\(\{0,\ldots,728\}\). Si \(b_{n+1}\) est le bloc admissible suivant,
le prolongement positionnel vérifie
\[
 N_{n+1}=729N_n+\nu(b_{n+1}),\qquad
 a_n=\frac{N_n}{729^n},\qquad b_n^+=\frac{N_n+1}{729^n}.
\]
Ici, \(N_0\) est la partie entière que détermine le lecteur régional. La lettre \(b_n^+\) désigne une extrémité réelle et se distingue du bloc ternaire \(b\). Le cylindre de représentation est \([a_n,b_n^+)\) ; son adhérence est utilisée dans l'argument d'existence de la limite. La retenue conserve le passage d'un préfixe au précédent par division euclidienne.

\begin{proposition}[Limite des centres cylindriques]
Pour un prolongement compatible, les adhérences des cylindres sont
emboîtées et déterminent une unique coordonnée \(x\). Si
\(c_n=(a_n+b_n^+)/2\), alors
\[
 x=\lim_{n\to\infty}a_n=\lim_{n\to\infty}b_n^+
 =\lim_{n\to\infty}c_n,\qquad
 |x-c_n|\leq\frac1{2\,729^n}.
\]
\end{proposition}
\begin{proof}
L'inégalité \(0\leq\nu(b_{n+1})<729\) situe le cylindre successeur dans l'adhérence du prédécesseur. Ses extrémités sont rationnelles, sa largeur est \(729^{-n}\) et elles demeurent dans le premier intervalle fermé borné. L'emboîtement et la largeur décroissante produisent une unique classe de Cauchy d'extrémités. Dans la réalisation archimédienne, cette classe a pour représentant \(x\) ; la distance au point milieu est au plus égale à la moitié de la largeur. La convention de frontière attribue une écriture lorsque deux développements positionnels représentent la même limite.
\end{proof}

La moyenne arithmétique intervient donc dans une évaluation précise : le
centre de chaque intervalle qui borne la coordonnée. La limite conserve
l’exactitude de la valeur bien que la largeur disparaisse. La représentation
enrichie conserve en outre la suite des préfixes et leur incidence.
La projection numérique et la projection exceptionnelle sont deux applications
à codomaines différents ; l’une conserve la valeur limite et l’autre transporte
les relations qui permettent la construction réticulaire. Le passage à
\(\RR\) se situe ainsi dans la réalisation du lecteur, après la
génération arithmétique des données qu’il évalue.

Pour expliciter la compatibilité, soient \(p_{n+1,n}\) les restrictions de l'état enrichi et \(x_n=p_{n+1,n}(x_{n+1})\). Si \(I_n(x_n)\) est le cylindre numérique et \(\mathcal B_n(x_n)\) l'incidence de frontière, les deux transports conservent simultanément
\[
 \overline{I_{n+1}(x_{n+1})}\subseteq\overline{I_n(x_n)},\qquad
 r_{n+1,n}\mathcal B_{n+1}(x_{n+1})=\mathcal B_n(x_n).
\]
La seconde égalité utilise les applications d'incidence du corpus ; la proposition~\ref{prop:lf-incidence-naturality} explicite leur réalisation au moyen de repères et de codes transportés. Ses cinq constructions consubstantielles ---spectrale, solénoïdale, de jauge, cohomologique et déterminantale--- appartiennent à ce même système compatible. L'étude géométrique ultérieure conserve cette unité d'origine en spécifiant laquelle de ses observations elle utilise dans chaque théorème.

\subsection{Trajectoires orientées et temporalité de l’actualisation}

Soit \(x\xrightarrow{\rm adm}y\) la relation d’admissibilité du TPK dans
l’espace enrichi \(\widehat X\). Une trajectoire paramétrée par
un ensemble discrètement ordonné \(\mathcal T\) est
\[
 \gamma:\mathcal T\longrightarrow\widehat X,
 \qquad\gamma(t)\xrightarrow{\rm adm}\gamma(t^+).
\]
La relation de compatibilité détermine les transitions ; le paramétrage
permet de les ordonner. Dans l’espace \(E=\mathcal T\times\widehat X\),
avec \(p(t,x)=t\), la trajectoire détermine la section
\(\widetilde\gamma(t)=(t,\gamma(t))\), qui vérifie
\(p\circ\widetilde\gamma=\id_{\mathcal T}\).
Cette précision conserve la différence entre un état, sa trajectoire et
la coordonnée utilisée pour décrire l’évolution.

La signature tritique \(\tau\in\{+1,0,-1\}\) type les générateurs au moyen de
\[
 J_\tau^2=-\tau I.
\]
Dans le lecteur temporel du corpus, \(+1\) représente le retour et la mémoire ; \(0\), la mise à jour de seuil ; et \(-1\), l'ouverture de prolongements conjugués. Leur interprétation comme passé, présent et futur se rapporte aux fonctions de l'antécédent conservé, de la transition effective et de la continuation admissible. Les trois temps verbaux acquièrent ainsi un contenu opératoriel sur les histoires.

Si \(\mathcal P_n(x)\) est l'ensemble des continuations admissibles de longueur \(n\), la suppression de la dernière étape définit \(r_n:\mathcal P_{n+1}(x)\to\mathcal P_n(x)\). Une continuation illimitée est une collection \((\gamma_n)\) avec \(r_n(\gamma_{n+1})=\gamma_n\). La restriction reconstruit les antécédents ; l'actualisation \(U_t\) agit dans la transition ; le prolongement maintient la frontière du préfixe. L'existence de ces collections est utilisée dans les domaines survivants fixés par le théorème de prolongement du corpus.

Sur une signature finie \(\sigma=(\tau_1,\ldots,\tau_n)\), l’inversion
orientée est
\[
 C_{\rm rev}\sigma=(-\tau_n,\ldots,-\tau_1),
 \qquad C_{\rm rev}^2\sigma=\sigma.
\]
La trajectoire conjuguée conserve également les transformations de feuillet, de position et de frontière de son transport. La bidirectionnalité s'exprime donc dans l'historique complet, tandis que cette formule exprime sa signature. Le retour de phase de \(\Gamma_9\) admet un incrément de mémoire : des états de phase identique peuvent occuper des profondeurs différentes.

\subsection{Algèbres quadratiques et activité du seuil parabolique}

Dans un régime fixe, la conjugaison de \(aI+bJ_\tau\) change le signe de \(b\), et sa norme algébrique est
\[
 (aI+bJ_\tau)(aI-bJ_\tau)=(a^2+\tau b^2)I.
\]
Les types elliptique, parabolique et hyperbolique correspondent aux trois
relations quadratiques. Leurs groupes locaux sont représentés par
\[
 e^{tJ_{+1}}=\cos t\,I+\sin t\,J_{+1},\quad
 e^{tJ_0}=I+tJ_0,\quad
 e^{tJ_{-1}}=\cosh t\,I+\sinh t\,J_{-1}.
\]
La preuve sépare les termes pairs et impairs de la série exponentielle ; dans le type nilpotent, seuls demeurent les degrés zéro et un. Chaque collection satisfait la loi de composition dans son paramètre et l'inversion \(e^{-tJ_\tau}=(e^{tJ_\tau})^{-1}\).

La valeur tritique zéro permet en particulier un générateur distinct de
zéro dont le carré est nul. Pour
\(N=\left(\begin{smallmatrix}0&1\\0&0\end{smallmatrix}\right)\),
\[
 e^{tN}(q,p)=(q+tp,p).
\]
Le seuil possède une action linéaire effective sur sa fibre. Son paramètre,
son déplacement et la mémoire de la transition restent déterminés
même lorsque la signature de régime est nulle. Tel est le contenu précis qui
permet de développer dans l’exposé le présent comme mise à jour active.

La forme elliptique positive a des niveaux circulaires en coordonnées
orthonormées et elliptiques après une transformation linéaire ; la forme hyperbolique
possède deux directions isotropes et des niveaux hyperboliques ; le type parabolique
possède une forme dégénérée et un transport unipotent. Une géométrie sphérique
requiert en outre l’espace homogène ou la métrique qui la réalise. La
classification précédente fixe les générateurs qui peuvent intervenir dans cette
réalisation et maintient séparés le type algébrique et la courbure d’un
espace ultérieur.

Le calendrier des modes concrétise la connexion avec l'auto-échelle. La règle \(+1:m\mapsto\Sigma\), \(-1:m\mapsto\Pi\), \(0:m\mapsto m\) appliquée au bloc cyclique \((+1,-1,0)\) produit \((\Sigma,\Pi,\Pi)\). Compter ses trois transitions, dans l'ordre \((\Sigma,\Pi)\), détermine
\begin{equation}
 F_{\rm av}=\begin{pmatrix}0&1\\1&1\end{pmatrix}.
 \label{eq:geo-fav}
\end{equation}
Son rayon propre positif, normalisé sous la forme \((1,x)^\mathsf T\), satisfait \(x=\lambda\), \(1+x=\lambda x\), et donc \(x^2-x-1=0\). La racine positive est la coordonnée d'auto-échelle exprimée par \(\varphi\). La matrice provient du calendrier des opérations ; l'expression radicale et les identités exponentielles suivantes sont ses réalisations ultérieures.

% Articulación posterior a la generación de pi, phi y e.
% Contenido acumulativo; no reemplaza el desarrollo previo del TRIT.
\subsection{Réalisations exponentielles des coordonnées générées}
\label{geo-euler-trit-articulacion}

Les coordonnées régionales déjà déterminées permettent de reconnaître les
réalisations concrètes de la collection quadratique du TRIT. L'ordre est
essentiel : la construction des régions précède l'évaluation d'une
exponentielle aux paramètres \(\pi\) ou \(2\log\varphi\). Ces expressions
identifient la fonction des coordonnées dans des opérateurs construits à partir
d'APP et du TPK ; elles n'interviennent pas dans la sélection des préfixes qui les produisent.

\subsubsection{Cycle ternaire, transport nilpotent et incidence des feuillets}

La multiplication \(a(r)=4r\) dans \(\mathbb Z/9\mathbb Z\) est d’ordre
\(3\). Sur les unités, elle détermine les orbites \(\{1,4,7\}\) et
\(\{2,8,5\}\). Dans le module d’augmentation de l’une quelconque d’entre elles, formé
des combinaisons entières dont les coefficients ont une somme nulle, la base
\((e_r-e_{4r},e_{4r}-e_{7r})\) donne
\[
 C=\begin{pmatrix}0&-1\\1&-1\end{pmatrix},
 \qquad C^2+C+I=0.
\]
La forme restreinte a pour matrice de Gram
\(\bigl(\begin{smallmatrix}2&-1\\-1&2\end{smallmatrix}\bigr)\).
Dans les évaluations résiduelles modulo \(9\), l’action distingue les feuillets :
\(S(ax,ay)=aS(x,y)\), tandis que \(P(ax,ay)=a^{-1}P(x,y)\).
La réalisation cyclique conserve donc une orientation contragrédiente
entre somme et produit. Les quotients du relèvement entier sont conservés
dans le registre arithmétique correspondant.

Le transport parabolique élémentaire est représenté par
\[
 N=\begin{pmatrix}0&1\\0&0\end{pmatrix}.
\]
L’incidence surfacique--volumique déjà construite dans \eqref{eq:geo-fav} fournit
\[
 F_{\rm av}=\begin{pmatrix}0&1\\1&1\end{pmatrix},
 \qquad A=F_{\rm av}^2=\begin{pmatrix}1&1\\1&2\end{pmatrix}.
\]
Les trois opérateurs agissent dans des espaces réels déclarés : le module
d’augmentation tensorisé avec \(\mathbb R\), le plan du transport nilpotent et
le plan d’incidence modale. Le partage d’une dimension matricielle n’identifie
ni ces espaces ni leurs coordonnées.

\begin{proposition}[Générateurs concrets des trois types]
\label{geo-euler-trit-generadores}
Les opérateurs
\[
 J_{\mathrm e}=\frac{2C+I}{\sqrt3},\qquad
 J_{\mathrm p}=N,\qquad
 J_{\mathrm h}=\frac{2A-3I}{\sqrt5}
\]
satisfont, respectivement,
\[
 J_{\mathrm e}^{\,2}=-I,\qquad
 J_{\mathrm p}^{\,2}=0,\qquad
 J_{\mathrm h}^{\,2}=I.
\]
\end{proposition}
\begin{proof}
La relation de \(C\) implique \((2C+I)^2=-3I\).
La multiplication directe donne \(N^2=0\).
Enfin, \(A\) est de trace \(3\) et de déterminant \(1\), d’où
\(A^2-3A+I=0\) et \((2A-3I)^2=5I\).
\end{proof}

\begin{proposition}[Trois évaluations de l’identité exponentielle]
\label{geo-euler-trit-evaluaciones}
Dans les réalisations précédentes,
\begin{equation}
 \boxed{
 \exp(\pi J_{\mathrm e})+I=0,\qquad
 \exp(J_{\mathrm p})=I+J_{\mathrm p},\qquad
 \exp(2\log\varphi\,J_{\mathrm h})=F_{\rm av}^{\,2}.}
 \label{geo-euler-trit-tres-evaluaciones}
\end{equation}
\end{proposition}
\begin{proof}
La spécialisation elliptique de l’exponentielle tritique donne
\(\cos\pi\,I+\sin\pi\,J_{\mathrm e}=-I\).
La nilpotence élimine exactement les termes de degré au moins \(2\)
dans la deuxième expression. Pour la troisième, \(\varphi^2=\varphi+1\) implique
\[
 \cosh(2\log\varphi)=\frac32,\qquad
 \sinh(2\log\varphi)=\frac{\sqrt5}{2}.
\]
L’exponentielle vaut donc
\(\frac32I+\frac{\sqrt5}{2}J_{\mathrm h}=A\).
\end{proof}

La première égalité représente le demi-tour elliptique ; le tour complet
satisfait \(\exp(2\pi J_{\mathrm e})=I\). La deuxième maintient un terme
linéaire non nul : le régime parabolique exprime un transport, et non une absence
d'évolution. La troisième réalise une avancée hyperbolique unimodulaire. Ses
valeurs propres sont \(\varphi^2\) et \(\varphi^{-2}\) ; une direction se dilate et
l'autre se contracte en conservant le déterminant. Les trois évaluations
emploient des paramètres distincts d'une collection exponentielle commune.

Dans cette collection, \(e\) intervient par l'évaluation scalaire
\(\exp(I)=eI\) et la loi \(\exp(sI)\exp(tI)=\exp((s+t)I)\).
Son rôle ne se limite pas au type parabolique. La coordonnée \(e\), obtenue
auparavant à partir de sa région, est reconnue ici comme l'évaluation en l'unité
de l'exponentielle ; cette identité ne remplace pas sa généalogie coinductive.

\subsubsection{Résolution polynomiale des régimes}

Les trois réalisations admettent une présentation compacte sur la somme
directe de leurs espaces. On définit
\[
 \mathbb J=J_{\mathrm e}\oplus J_{\mathrm p}\oplus J_{\mathrm h},
 \qquad Q=\mathbb J^2.
\]
Les relations précédentes impliquent
\[
 \mathbb J^6=\mathbb J^2,\qquad Q^3=Q.
\]
L’emploi de \(Q\) maintient ce carré opératoriel distinct du registre
dodécaphasique \(K\).

\begin{proposition}[Projecteurs polynomiaux de régime]
\label{geo-euler-trit-proyectores}
Les applications
\[
 p_{\mathrm e}=\frac{Q^2-Q}{2},\qquad
 p_{\mathrm p}=I-Q^2,\qquad
 p_{\mathrm h}=\frac{Q^2+Q}{2}
\]
sont des idempotents mutuellement orthogonaux, de somme \(I\), qui restituent les trois
termes de la somme directe représentant les régimes.
\end{proposition}
\begin{proof}
Les trois polynômes prennent les valeurs \((1,0,0)\), \((0,1,0)\) et
\((0,0,1)\) aux points \(-1,0,+1\), respectivement. Les différences
entre leurs carrés et eux-mêmes, ainsi que leurs produits croisés, sont
des multiples de \(X^3-X\). Substituer \(Q\), qui annule ce polynôme, démontre
les identités. Dans les trois blocs, \(Q\) agit comme \(-I,0,I\).
\end{proof}

On obtient ainsi
\begin{align}
 \exp(t\mathbb J)
 ={}&p_{\mathrm e}(\cos t\,I+\sin t\,\mathbb J)
   +p_{\mathrm p}(I+t\mathbb J)\notag\\
   &+p_{\mathrm h}(\cosh t\,I+\sinh t\,\mathbb J).
 \label{geo-euler-trit-exponencial-total}
\end{align}
Chaque terme conserve sa relation quadratique et la conjugaison inverse
son évolution. La représentation totale réunit les trois régimes, mais ne
définit pas à elle seule un opérateur de transition entre eux. Une telle transition
requiert les applications de transport du TPK avec leurs domaines, leur orientation et
leur mémoire. L’articulation exponentielle permet de reconnaître leurs types sans
remplacer cette dynamique par une identification des fibres.

% PROCEDENCIA FOCAL (no impresa)
% Resultado recuperado: tres generadores y tres evaluaciones exponenciales.
% Propietario completo:
% BASE/manuscrito/sections/hmt/09b_algebras_cuadraticas.tex:289-516.
% Antecedente causal de C y acción contragrediente:
% BASE/manuscrito/sections/hmt/03d_esqueleto_ciclotomico_app.tex:13-47,145-246.
% Los tres tipos y la orientación están en:
% BASE/manuscrito/sections/hmt/02_trit_geometrias.tex:83-196.
% Resolución polinómica recuperada:
% output/INVESTIGACION_HOLONOMIA_NONADICA_CRISTAL_APERIODICO_20260903/
% ALGEBRA_HOLONOMICA_UNIVERSAL_Y_EULER_TRITICO.md, secciones 8-11 y 22.
% Los cuatro propietarios se leyeron completos. BASE designa:
% /Users/ruben/Documents/HMT2/
% EL_CIERRE_HOLOGRAFICO_DEL_INFINITO_SUCESOR_102_CAPITULOS_2026-08-28_EN_TRABAJO/
% No se incorpora la identidad P9(alpha)=delta del documento especializado:
% esa identificación no pertenece a este bloque.
% Tampoco se identifica una dirección nilpotente con la vacancia electrónica,
% ni se presume una representación global del transporte por la suma directa.
% COMPROBACIONES: Python estándar, enteros y Fraction.
% C^2+C+I=0; (2C+I)^2=-3I; N^2=0; (2Fav^2-3I)^2=5I.
% 81 pares de residuos verifican las dos acciones contragredientes.
% Los 3 proyectores se verificaron en Q=-1,0,+1; prueba general en texto.
% 6 labels únicos y entornos equilibrados. Sin compilación.

